\chapter[Universal Geometric Bounds]{Universal Geometric Bounds}\label{ch:sec:universal-geometric-bounds}

The two universal theorems of this chapter---the Lipschitz bound and the quadratic bound---apply to every sufficiently regular functional on the WIS manifold. They are the technical core of Part~IV and will be specialised to Bregman divergences in Chapter~\ref{ch:sec:the-average-hessian-operator-and-universal-bregman-theo} and to piecewise-smooth functionals in Chapters~17--18.


\section{Introduction}\label{ch:sec:introduction-universal_bounds}

We have seen that the UOD bounds the variation of individual entropies.
But the pattern is always the same: $|S(P)-S(Q)|\le L\sqrt{\mathcal
D(P,Q)}$.  This is not a coincidence.  The inequality follows from the
geometry of the manifold, not from the specific functional.  This
chapter extracts the common geometric core: two universal theorems that
apply to \emph{every} sufficiently regular functional on the WIS
manifold, establishing Lipschitz and quadratic bounds purely in terms
of the UOD.
Throughout this chapter, let $(\mathcal P,g)$ denote the posterior
manifold constructed in Part~\ref{ch:part:posterior}, let
$F:\mathcal P\rightarrow\mathbb R^{n}$ be the global WIS coordinate map,
and let $\mathcal D(P,Q)=\|F(P)-F(Q)\|^{2}$ be the two-point Universal
Optimality Defect.  Under the identification established in Part~\ref{ch:part:posterior},
this coincides with the squared geodesic distance,
$\mathcal D(P,Q)=d_g(P,Q)^{2}$.
\paragraph{Smooth Functionals.}
Let \(S:\mathcal P\rightarrow\mathbb R\) be a continuously differentiable
functional.  Its Riemannian gradient is denoted \(\nabla_g S\).
We assume \(\|\nabla_g S(P)\|_g\le L\) for every point of the region under consideration.

\begin{theorem}[Universal Lipschitz Bound]\label{ch:thm:14-1}

Let \(S\in C^{1}(\mathcal P)\) satisfy
\(\|\nabla_gS\|_g\le L.\)
Then, for every pair of posterior distributions \(P,Q\),
\(\|S(P)-S(Q)\| \le L\,d_g(P,Q).\)
Consequently,
\[ \boxed{
|S(P)-S(Q)|
\le
L\sqrt{\mathcal D(P,Q)}.
} \]
\end{theorem}

\begin{quote}
\textbf{Mathematical intuition.}
The Lipschitz bound $|S(P)-S(Q)|\le L\sqrt{\mathcal D(P,Q)}$ says that functionals cannot vary faster than the square root of the UOD.
Since $\sqrt{\mathcal D}$ is the geodesic distance on the manifold, this is exactly the definition of Lipschitz continuity with respect to the manifold metric.
\end{quote}

\begin{proof}
Let \(\gamma(t)\) be a minimizing geodesic joining \(P\) to \(Q\) in the
pullback metric, parametrized so that \(\gamma(0)=P\) and \(\gamma(1)=Q\).
By the fundamental theorem of calculus,
\[
S(Q)-S(P)=\int_0^1 \langle\nabla_g S(\gamma(t)),\dot\gamma(t)\rangle_g\,dt.
\]
Applying the Cauchy--Schwarz inequality and the bound on the gradient,
\[
|S(Q)-S(P)|
\le \int_0^1 \|\nabla_g S(\gamma(t))\|_g\,\|\dot\gamma(t)\|_g\,dt
\le L\int_0^1 \|\dot\gamma(t)\|_g\,dt
= L\,d_g(P,Q).
\]
Since \(d_g(P,Q)=\sqrt{\mathcal D(P,Q)}\) by construction, the
claimed inequality follows.
\end{proof}

\section{Local Quadratic Bound}\label{ch:sec:local-quadratic-bound}
The previous theorem only requires first-order regularity.
Second-order information yields stronger local estimates.
\begin{theorem}[Local Quadratic Bound]
\label{ch:thm:14-2}
Let
$S\in C^{2}(\mathcal P),$
and suppose
$\nabla_gS(P^{\star})=0.$
Then there exists a neighborhood of \(P^{\star}\) and a constant
\(C>0\) such that
\[ \boxed{
|S(P)-S(P^\star)|
\le
C\,\mathcal D(P,P^\star).
} \]
\end{theorem}
\begin{proof}[Sketch]
Let \(\gamma(t)\) be the geodesic joining \(P^{\star}\) to \(P\) in the
pullback metric, parametrized so that \(\gamma(0)=P^{\star}\) and
\(\gamma(1)=P\).  Taylor's theorem applied to the composition
\(S\circ\gamma\) gives
\[
S(P)=S(P^{\star})+\langle\nabla_g S(P^{\star}),\dot\gamma(0)\rangle_g
+\frac12\,\ddot\gamma(0)^T\operatorname{Hess}_g S(P^{\star})\dot\gamma(0)
+O\bigl(\mathcal D(P,P^{\star})^{3/2}\bigr).
\]
Since \(\nabla_g S(P^{\star})=0\) by hypothesis, the linear term vanishes.
The quadratic term is bounded by \(\frac12\|\operatorname{Hess}_g S\|\,
\mathcal D(P,P^{\star})\).  The estimate follows from boundedness of the
Hessian on the compact region under consideration. \(\)
\end{proof}
\paragraph{Interpretation.}
The previous two theorems establish two universal regimes.
First-order regularity produces
\(O(\sqrt{\mathcal D}),\)
while second-order regularity around stationary points improves the
estimate to
\(O(\mathcal D).\)
These results depend only on the geometry of the WIS manifold and not on
the specific functional.
\section{Consequences}\label{ch:sec:consequences-universal_bounds}
Any information functional satisfying the hypotheses of Theorem~\ref{ch:thm:14-1}
immediately inherits a geometric stability estimate.
Likewise, every twice continuously differentiable functional having a
stationary point satisfies the quadratic estimate of Theorem~\ref{ch:thm:14-2}.
Therefore no functional-specific argument is required to establish
continuity or local stability.
Only verification of the regularity assumptions remains.
\paragraph{Transition to Chapter~\ref{ch:sec:universal-geometric-bounds}.}
\paragraph{Running example: differential privacy.}
In differential privacy (Section~\ref{ch:sec:differential-privacy-and-database-queries}), the Laplace mechanism produces a smooth privacy loss $\mathcal L$.  Theorem~\ref{ch:thm:14-1} applies with Lipschitz constant $L\le 1/\min_i\pi_i$, providing an explicit uniform bound on how much the privacy loss can change across the posterior manifold.  The geo-indistinguishability example (Section~\ref{ch:sec:geo-indistinguishability-and-location-privacy}) satisfies the same hypotheses, and Theorem~\ref{ch:thm:14-2} sharpens the estimate near optimal posteriors.
The results of this chapter apply to arbitrary smooth scalar fields.
The next chapter specializes these universal estimates to Bregman
divergences, introducing curvature operators and a geometric comparison
theory that encompasses a broad class of information-theoretic
divergences within the WIS framework.
